% ---------------------------------------------------------------------------
% Abstract: two versions (approx. 300 and 500 words).
% Choose which is used with the switch in main.tex:
%   \shortabstracttrue   -> ~300-word abstract
%   \shortabstractfalse  -> ~500-word abstract
% ---------------------------------------------------------------------------

\section*{Abstract}

\ifshortabstract
% ======================= ~300-word abstract =======================

\textbf{Background:} Intravenous thrombolysis (IVT) and mechanical thrombectomy (MT) are both time-dependent, but randomised trials have largely studied each in isolation. In routine care, many patients receive IVT at one hospital and MT, sometimes hours later, at another. We estimated the real-world effectiveness of IVT and MT, alone and combined, and how it depends on the timing of each.

\textbf{Methods:} We used national stroke registry data for 74,678 patients with ischaemic stroke and an NIHSS score above 5 (or a recorded large vessel occlusion) who arrived within six hours of onset at MT-performing hospitals in England and Wales, 2020--2025. A good outcome was a discharge modified Rankin Scale (mRS) score of 0--2, or no worse than before the stroke. We used three causal methods with different assumptions: target trial emulation by matching, an explainable S-learner, and an orthogonal causal forest. Treatment was classed as early ($\leq$270 minutes from onset) or late.

\textbf{Results:} 28,884 patients received IVT alone, 2,983 MT alone and 5,164 both. Every strategy improved outcomes, and all three methods ranked them identically. Early combined treatment gave the largest benefit: 46\% of treated patients versus 24\% of matched controls had a good outcome (causal forest: about 18 percentage points). Early MT alone and early IVT alone gave gains of about 14 and 10 percentage points. Benefit declined approximately linearly with delay, reaching zero at about 6.8 hours for IVT and 9.9 hours for MT; each hour of delay multiplied the odds ratio by about 0.84 (IVT) and 0.86 (MT). Early IVT added most when MT was delayed.

\textbf{Conclusions:} In routine practice, IVT and MT are each effective, with time dependence closely matching the trials. The greatest benefit comes from giving both treatments early; stroke services should aim to minimise both treatment times, not either alone.

\else
% ======================= ~500-word abstract =======================

\textbf{Background:} Intravenous thrombolysis (IVT) and mechanical thrombectomy (MT) are highly effective stroke treatments, but the benefit of each declines with delay. Trials have largely studied each treatment and its timing in isolation, and trials of MT with or without prior IVT recruited patients in whom the two treatments were given close together. In routine care, many patients receive IVT locally and are transferred for MT, which may follow hours later. We estimated the real-world effectiveness of IVT and MT, alone and combined, and how it depends on the time to each treatment, separately and together.

\textbf{Methods:} We used data from the Sentinel Stroke National Audit Programme (SSNAP) for 74,678 patients with ischaemic stroke and an NIHSS score above 5 (or a recorded large vessel occlusion) who arrived within six hours of known onset at MT-performing hospitals in England and Wales between 2020 and 2025. A good outcome was a discharge modified Rankin Scale (mRS) score of 0--2, or no worse than before the stroke. Treatment was classed as early ($\leq$270 minutes from onset) or late ($>$270 minutes). We used three causal inference methods with different assumptions: target trial emulation by matching; an explainable S-learner (XGBoost), predicting each patient's outcome with and without treatment; and an orthogonal, doubly robust causal forest. Effects were estimated in treated patients. A secondary model with SHAP explanations identified the characteristics associated with greater benefit.

\textbf{Results:} Overall, 28,884 patients received IVT alone, 2,983 MT alone, 5,164 both, and 37,647 neither. Every strategy was associated with a higher chance of a good outcome, and all three methods ranked the strategies identically. Early combined treatment gave the largest benefit: 46\% of treated patients against 24\% of matched controls had a good outcome, the S-learner estimated an odds ratio of about 4.4, and the causal forest an absolute gain of about 18 percentage points. Early MT alone and early IVT alone gave absolute gains of about 14 and 10 percentage points. Late treatments gave roughly half to two thirds of the benefit, on the log odds scale, of their early counterparts, and patients receiving early IVT with late MT gained little more than those receiving early IVT alone. With treatment times modelled continuously, benefit declined approximately linearly, reaching zero at about 6.8 hours for IVT and 9.9 hours for MT; each hour of delay multiplied the odds ratio by about 0.84 (IVT) and 0.86 (MT). Joint analysis showed that benefit was greater the earlier IVT was given, for any time to MT, and the earlier MT followed, for any time to IVT. The admitting stroke team, stroke severity and pre-stroke disability were most associated with the size of benefit.

\textbf{Conclusions:} In routine practice, IVT and MT are each effective, with benefits and time dependence closely matching the randomised trials. Early IVT adds substantially to MT, especially when MT is delayed, whereas MT adds little to early IVT once delayed beyond about six hours. The greatest benefit comes from giving both treatments early, and services should minimise both treatment times, not either alone.

\fi
